% 05-preface —— 由 tools/convert.py 自动生成，勿手改
% 源文件：part0000_split_007.html, part0000_split_008.html, part0000_split_009.html, part0000_split_010.html, part0000_split_011.html, part0000_split_012.html, part0000_split_014.html, part0000_split_015.html

Rust 是一门用于系统编程的语言。

事实上，大多数程序员并不熟悉系统编程，尽管它是我们日常工作的基础。

假设你合上了笔记本计算机，操作系统检测到这个行为，随即暂停了所有正在运行的程序、关闭屏幕，并让计算机进入睡眠状态。稍后，当你又打开笔记本计算机时，屏幕和其他组件再次启动，每个程序都能从断点处恢复运行。这些我们习以为常的功能，离不开系统程序员编写的大量系统程序。

系统编程用于：

\begin{itemize}
\item 操作系统
\item 各种设备驱动程序
\item 文件系统
\item 数据库
\item 在成本极低或可靠性要求极高的设备上运行的代码
\item 加密解密
\item 媒体编解码器（用于读取和写入音频、视频和图像文件的软件）
\item 媒体处理（例如，语音识别或照片编辑软件）
\item 内存管理（例如，实现垃圾回收器）
\item 文本渲染（将文本和字体转换为像素）
\item 实现高级编程语言（如 JavaScript 和 Python）
\item 网络编程
\item 虚拟化和软件容器
\item 科学仿真
\item 游戏
\end{itemize}

简而言之，系统编程是\textbf{资源受限}条件下的编程。每字节、每个 CPU 周期，对于程序正常运行都尤为重要。

即便是支持一个最基本的应用程序，涉及的系统代码量也是惊人的。

本书并不会教你如何进行系统编程。虽然本书涵盖了内存管理的许多细节，但是如果你还没有亲自做过系统编程，那么这些细节乍一看可能会有点儿过于深奥。如果你是一名经验丰富的系统程序员，就会看出 Rust 的卓尔不群：这种新工具可以解决困扰整个行业数十年的、众所周知的重大问题。

\section{读者对象}

如果你已经是系统程序员并且正准备物色一款 C++ 的替代品，那么本书适合你。如果你是任何一种编程语言的资深开发人员，无论是 C\#、Java、Python、JavaScript，还是其他语言，那么本书同样适合你。

不过，你不能仅仅学习 Rust 语言本身。要充分发挥该语言的价值，还需要接触一些系统编程知识。我们建议你在阅读本书的同时，使用 Rust 实现一些系统编程领域的业余项目，利用 Rust 的速度、并发和安全，构建出你以前从未构建过的东西。开头那个主题列表应该会给你一些启发。

\section{写作初衷}

最初学习 Rust 时，我们曾期待有这样一本书，这正是本书的写作初衷。我们的目标是直面 Rust 中一些重大而新颖的概念，清晰而深入地呈现它们，从而减少学习中的试错成本。

\section{浏览本书}

本书前两章简要介绍了 Rust 的背景知识，接着在第 3 章开始介绍基本数据类型。第 4 章和第 5 章讨论了“所有权”和“引用”这两个核心概念。建议你按顺序通读前 5 章。

第 6 章到第 10 章涵盖了 Rust 这门语言的基础知识：表达式（第 6 章）、错误处理（第 7 章）、crate 与模块（第 8 章）、结构体（第 9 章），以及枚举与模式（第 10 章）。这几章可以稍微读快一些，但不要跳过第 7 章。

第 11 章介绍了特型与泛型，这是最后两个你需要了解的重要概念。特型就像 Java 或 C\# 中的接口。它们也是 Rust 用来将你的类型集成到语言本身的主要方式。第 12 章展示了如何用特型支持运算符重载，第 13 章介绍了更多的实用工具特型。

了解特型和泛型可以帮你解锁本书其余部分的知识。不容错过的两个强大工具——闭包和迭代器，会分别在第 14 章和第 15 章中进行介绍。剩下的几章你可以按任意顺序阅读，或根据需要深入阅读。它们涵盖了这门语言的其余部分：集合（第 16 章）、字符串与文本（第 17 章）、输入与输出（第 18 章）、并发（第 19 章）、异步编程（第 20 章）、宏（第 21 章）、不安全代码（第 22 章），以及调用来自其他语言的函数（第 23 章）。

\section{排版约定}

本书使用以下排版约定。

\textbf{黑体}

表示新术语或重点强调的内容。

等宽字体（\texttt{constant width}）

表示程序片段，以及正文中出现的变量名、函数名、数据库、数据类型、环境变量、语句和关键字等。

\begin{mySidenote}
加粗等宽字体（\textbf{\texttt{constant width bold}}）
\end{mySidenote}

表示应该由用户输入的命令或其他文本。

\begin{mySidenote}
等宽斜体（\emph{\texttt{constant width italic}}）
\end{mySidenote}

表示应该由用户输入的值或根据上下文确定的值替换的文本。

此图标表示一般性注释。

\section{使用代码示例}

补充材料（代码示例、练习等）可在 \href{https://github.com/ProgrammingRust}{https://github.com/\allowbreak{}ProgrammingRust} 下载。

本书旨在帮助你完成工作。一般来说，你可以在自己的程序或文档中使用本书提供的示例代码。除非需要复制大量代码，否则无须联系我们获得许可。比如，使用本书中的几个代码片段编写程序无须获得许可，销售或分发 O'Reilly 图书的示例光盘则需要获得许可；引用本书中的示例代码回答问题无须获得许可，将本书中的大量示例代码放到你的产品文档中则需要获得许可。

我们很希望但并不强制要求你在引用本书内容时加上引用说明。引用说明一般包括书名、作者、出版社和 ISBN。比如“\emph{Programming Rust, Second Edition} by Jim Blandy, Jason Orendorff, and Leonora F. S. Tindall (O'Reilly ). Copyright 2021 Jim Blandy, Leonora F. S. Tindall, and Jason Orendorff, 978-1-492-05259-3”。

如果你对示例代码的用法超出了上述的许可范围，欢迎你通过 \href{mailto:permissions@oreilly.com}{permissions@oreilly.com} 与我们联系。

\section{联系我们}

如有与本书有关的评价或问题，请联系出版社。

美国：

O'Reilly Media, Inc.

1005 Gravenstein Highway North

Sebastopol, CA 95472

中国：

北京市西城区西直门南大街 2 号成铭大厦 C 座 807 室（100035）

奥莱利技术咨询（北京）有限公司

O'Reilly 的每一本书都有专属网页，你可以在那儿找到本书的相关信息，包括勘误表、示例代码\footnote{也可到图灵社区本书主页 \href{http://ituring.com.cn/book/2846}{ituring.com.cn/\allowbreak{}book/\allowbreak{}2846} 下载示例代码。——编者注}以及其他信息。本书的网页是 \href{https://oreil.ly/programming-rust-2e}{https://oreil.ly/\allowbreak{}programming-rust-2e}。

对于本书的评论和技术性问题，请发送电子邮件至 \href{mailto:errata@oreilly.com.cn}{errata@oreilly.com.cn}。

要了解更多 O'Reilly 图书、培训课程和新闻的信息，请访问以下网站：\texthref{https://www.oreilly.com}。

我们在 Facebook 的地址如下：\href{http://facebook.com/oreilly}{http://facebook.com/\allowbreak{}oreilly}。

请关注我们的 Twitter 动态：\href{http://twitter.com/oreillymedia}{http://twitter.com/\allowbreak{}oreillymedia}。

我们的 YouTube 视频地址如下：\href{http://www.youtube.com/oreillymedia}{http://www.youtube.com/\allowbreak{}oreillymedia}。

\section{致谢}

本书得以面世得益于我们的官方技术审稿人的关注，感谢 Brian Anderson、Matt Brubeck、J. David Eisenberg、Ryan Levick、Jack Moffitt、Carol Nichols 和 Erik Nordin。还要感谢本书各语种版本的翻译人员：Hidemoto Nakada（中田秀基）（日语）、李松峰（简体中文第 1 版）、汪志成（简体中文第 2 版），以及 Adam Bochenek 和 Krzysztof Sawka（波兰语）。

另外，许多非官方审稿人也阅读了早期的草稿并提供了宝贵的反馈。感谢 Eddy Bruel、Nick Fitzgerald、Graydon Hoare、Michael Kelly、Jeffrey Lim、Jakob Olesen、Gian-Carlo Pascutto、Larry Rabinowitz、Jaroslav Šnajdr、Joe Walker 和 Yoshua Wuyts 对本书给出的深思熟虑的评论。Jeff Walden 和 Nicolas Pierron 花大量时间几乎审阅了全书。与任何编程冒险一样，编程书总是因高质量的错误报告而日益兴盛。感谢你们。

Mozilla 非常支持吉姆和贾森在这个项目上的工作，尽管这不属于我们的官方职责范围，并且一定程度上会分散我们的注意力。非常感谢吉姆和贾森的经理 Dave Camp、Naveed Ihsanullah、Tom Tromey 和 Joe Walker 的支持。他们从长远的角度看待 Mozilla 的意义，我们希望这些成果能证明他们对我们的信任是对的。

还要感谢 O'Reilly 出版社帮助这个项目取得成果的每一个人，尤其是我们极富耐心的编辑 Jeff Bleiel 和 Brian MacDonald，以及我们的策划编辑 Zan McQuade。

最重要的是，衷心感谢我们的家人，感谢他们坚定不移的爱、热情和耐心。

